HAL Id: hal-01224644
https://hal.science/hal-01224644v1
Submitted on 16 Aug 2017
HAL is a multi-disciplinary open access archive
for the deposit and dissemination of scientific research documents, whether they are published or not.
The documents may come from teaching and research
institutions in France or abroad, or from public or private research centers.
L’archive ouverte pluridisciplinaire HAL, est destinée au dépôt et à la diffusion de documents scientifiques de niveau recherche, publiés ou non, émanant
des établissements d’enseignement et de recherche
français ou étrangers, des laboratoires publics ou
privés.
HAL Authorization
An Event-B formalization of KAOS goal refinement patterns
Abderrahman Matoussi, Frédéric Gervais, Régine Laleau
To cite this version:
Abderrahman Matoussi, Frédéric Gervais, Régine Laleau. An Event-B formalization of KAOS goal refinement
patterns. [Research Report] TR-LACL-2010-1, LACL. 2010. ⟨hal-01224644⟩
An Event-B formalization of KAOS goal refinement
patterns
Abderrahman Matoussi Fr´ed´eric Gervais R´egine Laleau
January 2010
TR–LACL–2010–1
Laboratoire d’Algorithmique, Complexit´e et Logique (LACL)
D´epartement d’Informatique
Universit´e Paris 12 – Val de Marne, Facult´e des Science et Technologie
61, Avenue du G´en´eral de Gaulle, 94010 Cr´eteil cedex France
Tel.: (33)(1) 45 17 16 47, Fax: (33)(1) 45 17 66 01
Laboratory of Algorithmics, Complexity and Logic (LACL)
University Paris 12 (Paris East)
Technical Report TR–LACL–2010–1
A. Matoussi, F. Gervais, R. Laleau.
An Event-B formalization of KAOS goal refinement patterns
c A. Matoussi, F. Gervais, R. Laleau, January 2010.
An Event-B formalization of KAOS goal refinement patterns
Abderrahman Matoussi Fr´ed´eric Gervais R´egine Laleau
Laboratory of Algorithmics, Complexity and Logic - University Paris 12 (Paris East),
France
{abderrahman.matoussi,frederic.gervais,laleau}@univ-paris12.fr
Abstract
Goals play an important role in requirements engineering process, and consequently
in systems development process. Whereas specifications allow us to answer the question
”WHAT the system does”, goals allow us to address the ”WHY, WHO, WHEN” questions [5]. Up to now, the main software development approaches using formal methods,
such as Event-B, begins at the specification level. Our objective is to include requirements analysis within this process, and more precisely the KAOS method, which is a
goal-oriented methodology for requirements engineering. The latter allows analysts to
build requirements models and to derive requirements documents. Existing work that
combine KAOS with formal methods generate a formal specification model from a KAOS
requirements model. We aim at expressing KAOS goal models with a formal language
(Event-B), hence staying at the same abstraction level. Thus we take advantage from the
Event-B method: (i) it is possible to use the method during the whole development process and (ii) we can benefit from the industrial maturity of tools supporting the method.
For that purpose, we propose a constructive approach in which Event-B models are built
incrementally from KAOS goal models, driven by goal refinement patterns.
Keywords. Requirements engineering, KAOS methodology, Event-B method, KAOS
refinement patterns.
1 Introduction
Employing formal methods for complex systems specification is steadily growing from year
to year. They have shown their ability to produce such systems for large industrial problems
such as Paris metro line 14 [6] or Roissy Val [7] using the B method [1]. With most of
formal methods, an initial mathematical model can be refined in multiple steps, until the final
refinement contains enough details for an implementation. Most of the time, the initial model
is derived from the description, obtained by the requirements analysis. Consequently, the
major remaining weakness in the development chain is the gap between textual or semi-formal
requirements and the initial formal specification. There is little research on reconciling the
requirements phase with the formal specification phase. In fact, the validation of this initial
formal specification is very difficult due to the inability to understand formal models (for
customers) and to link them with initial requirements (for designers).
Our objective is to combine the requirements and the specification phases by using KAOS
and the Event-B method. On one hand, KAOS is a goal-oriented methodology for requirements engineering which allows analysts to build requirements models and to derive require-
A. Matoussi, F. Gervais, R. Laleau. An Event-B formalization of KAOS goal refinement patterns 2
ments documents. On the other hand, Event-B is a model-based formal method which provides language, techniques and tools to support the analysis and design of systems, from the
specification to the implementation stages. Existing work [8, 9, 12, 13] that combine KAOS
with formal methods generate a formal specification model from a KAOS requirements model.
Contrary to these methods that take only a subset of the KAOS models into account, our
long-term objective is to express the whole KAOS requirements model with Event-B, in order
to support formal reasoning on it. The key idea is to stay at the same abstraction level as
KAOS.
In this report, we begin by considering the first stage of KAOS requirements analysis,
namely the goals modeling. In that aim, we present: i) an Event-B semantics for a subset
of the KAOS goals called ”Achieve” goals, and ii) a formalization in Event-B of the most
used refinement patterns of KAOS goals. This report continues our previous works [10, 11]
by addressing the Event-B formalization of KAOS pattern with additional studies, results
and proofs. The Event-B formalization of the other KAOS models is a work in progress.
The remainder of this report is organized as follows. Section 2 overviews the KAOS and the
Event-B formal methods that are employed in the proposed approach. Section 3 details our
proposed approach that consists to express KAOS goals model with Event-B. Sections 4, 5
and 6 present the Event-B formalization of, respectively, the milestone-driven goal refinement
pattern, the AND goal refinement pattern and the OR goal refinement pattern. Section 7
overviews some other KAOS goal refinement patterns. Section 8 illustrates the approach with
a case study. Relevant issues and related work are discussed in Section 9 and 10. Finally,
Section 11 concludes with an outline of future work.
2 Background
This section briefly describes the two methods that the proposed approach is based on namely
the Event-B formal method and the KAOS methodology.
2.1 KAOS method
KAOS (Knowledge Acquisition in autOmated Specification) [4] is a goal-based requirements
engineering method. KAOS requires the building of a data model in UML-like notation. A
goal defines an objective the system should meet, usually through the cooperation of multiple
agents such as devices or humans. KAOS differentiates between goals and domain properties
that are descriptive statements about the environment such as physical laws, organizational
norms or policies, etc. KAOS is composed of five complementary sub-models related through
inter-model consistency rules:
• The central model is the goal model which describes the goals of a system and its
environment. The core of the goal model consists of a refinement graph showing how
higher-level goals are refined (using the concept of refinement patterns [5]) into lowerlevel ones and, conversely, how lower-level goals contribute to higher-level ones. Higherlevel goals are strategic and coarse-grained while lower-level goals are technical and
fine-grained (more operational in nature).
• The object model defines the objects (agents, entity...) of interest in the application
domain.
A. Matoussi, F. Gervais, R. Laleau. An Event-B formalization of KAOS goal refinement patterns 3
• The agent responsibility model takes care of assigning goals to agents in a realizable
way.
• The operation model captures the system operations in terms of their individual
features and their links to the goal, object and agent models.
• The behavior model captures the required behaviors of system agents in terms of
temporal sequences of state transitions for the variables they control.
The importance of the goal model derives from the central role played by goals in the
requirements engineering (RE) process. For instance, we can derive other models in a systematic way from the goal model such as the object and operation models. Moreover, the
goal model enables early forms of RE-specific analysis such as risk analysis, conflict analysis,
or evaluation of alternative options.
KAOS provides a catalog of goal patterns that generalize the most common goal configurations. Achieve Goals specifies a property that the system will achieve “some time in
the future”. Cease Goals disallow achievement “some time in the future”. Maintain Goals
specifies a property that must hold “at all times in the future”. Avoid Goals prescribes a
property that must not hold “at all times in the future”.
Goals in KAOS can be either “AND” or “OR” refined. A goal is AND-refined into
subgoals, such that the conjunction of the subgoals is a sufficient condition to achieve the
parent goal. The OR-refinement associates a goal to a set of alternative subgoals in which the
achievement of the higher-level goal requires the achievement of at least one of its subgoals.
KAOS offers a lot of refinement patterns [5] that decompose goals. These patterns can only
be used in the context of different tactics defined in KAOS such as the milestone-driven
tactics which consists in identifying milestone states that must be reached to achieve the
target predicate.
KAOS also provides a criterion for stopping the refinement process. If a goal can be
assigned to the sole responsibility of an individual agent, there is no need for further goal
refinement to occur. Operational goals (goals that are assigned to agents) are the leaves of a
goal graph. Each leaf can be either a requirement (if it is assigned to an agent of the system)
or an expectation (if it is assigned to an agent in the environment). The reader may refer
to [4] for a full description of these notions.
Notice that KAOS provides an optional formal assertion layer for the specification of
goals in Real-Time Linear Temporal Logic (RT-LTL). This formalization step [5, 4] allows
to check for instance that goal refinements are correct and complete using a theorem prover,
formal refinement patterns, or a bounded SAT solver. Even if such checking is important
in order for example to detect missing subgoals in incomplete requirements, the use of this
kind of logic cannot fill in the gap between requirements and the later phases of development.
This is a serious shortcoming since it obliges designers to use another formal method for
developing their systems. Consequently, it is difficult to validate specifications with regard
to requirements even if they have been expressed with RT-LTL.
2.2 Event-B method
Event-B [3], an evolution of the classical B method [1], is a formal method for modeling
discrete systems by refinement. An Event-B model can be described in terms of two basic
constructs:
A. Matoussi, F. Gervais, R. Laleau. An Event-B formalization of KAOS goal refinement patterns 4
• The context: it provides axiomatic properties of Event-B models. It contains the
static part of a model such as carrier sets, constants, axioms and theorems. Carrier sets
are similar to types but both, carrier sets and constants, can be instantiated. Axioms
describe properties of carrier sets and constants. Theorems are derived properties
that can be proved from the axioms. Proof obligations associated with contexts are
straightforward: the stated theorems must be proved.
• The machine: it contains the dynamic part such as variables, invariants, theorems,
events and variants. Variables v define the state of a machine. Possible state changes
are described by means of events. Each event is composed of a guard G(t, v) and an
action S(t, v), where t are local variables the event may contain. The guard states the
necessary condition under which an event may occur, and the action describes how the
state variables evolve when the event occurs. The correctness of an Event-B model is
defined by an invariant property which every state in the system must satisfy. So, every
event in the system must be shown to preserve this invariant. In order to verify this
requirement, proof obligations have been defined.
It is also important to indicate that the most important feature provided by Event-B is its
ability to stepwise refine specifications. Refinement is a process that transforms an abstract
and non-deterministic specification into a concrete and deterministic system that preserves
the functionality of the original specification. During the refinement, event descriptions are
rewritten to take new variables into account. This is performed by strengthening their guards
and adding substitutions on the new variables. New events that only assign the new variables
may also be introduced. Proof obligations (POs) are generated to ensure the correctness of
the refinement with respect to the abstract model. Event-B is supported by several tools,
currently in the form a platform called Rodin [28].
3 Expressing Goals in Event-B
3.1 Motivation
In the software development life-cycle, the first phase corresponds to the requirements engineering. It is followed by the specification phase and then, the development stage. The
Event-B method has shown that it was very relevant for the last two phases. So, the approach
proposed in this report aims at introducing requirements analysis into the Event-B method.
Thus, it will be possible to establish formal links between this model and the specification of
a system. As we said before, we have chosen KAOS as a goal-oriented RE method because in
KAOS the emphasis was more on semi-formal and formal reasoning about behavioral goals
for derivation of goal refinements, goal operationalizations, etc, while in i* [20] for example
(the other famous goal-oriented RE), the emphasis was more on qualitative reasoning on soft
goals. The choice of Event-B is due to its similarity and complementarity with KAOS: (i)
both Event-B and KAOS have notions of refinement and constructive approach; (ii) KAOS
and Event-B (conversely to the classical B) have the ability to model both the system and its
environment. Hence, these points offer the possibility of using Event-B in the early phases
of the development of complex systems. This can help the designers in constructing a mathematical simulation of the overall system.
Since goals play an important role in requirements engineering process and provide a
bridge linking stakeholder requests to system specification [24], the proposed approach comes
A. Matoussi, F. Gervais, R. Laleau. An Event-B formalization of KAOS goal refinement patterns 5
down to systematically derive an Event-B representation from a KAOS goal model. Consequently, we show that it is possible to express KAOS goal models with formal method like
Event-B by staying at the same abstraction level. However, it is not possible to verify that
both models are equivalent since KAOS models are only semi-formal. The Event-B expression of a KAOS goal model allows to give it a precise semantics just as existing translations
from UML specifications into B specifications give a formal semantics to class diagrams or
state diagrams [18, 19].
3.2 Achieve Patterns
To achieve our objective, we formalize with Event-B the KAOS refinement patterns that
analysts use to generate a KAOS goal hierarchy. We think that these formal design patterns or
proof-based design patterns will be very useful and explores the fact that the Event-B method
provides a framework for developing generic models of systems. In this report, we focus on the
most frequently used goal patterns: the Achieve goals. Formalization of the other categories
of goal patterns is a work in progress. An Achieve goal prescribes intended behaviors where
some target condition must sooner or later hold whenever some other condition holds in the
current system state (this state is an arbitrary current one). An Achieve goal in KAOS is
denoted as follows: Achieve[TargetCondition From CurrentCondition]. This notation
has the following informal temporal pattern where CurrentCondition prefix is optional
(said otherwise, it can be true):
[if CurrentCondition then] sooner-or-later TargetCondition.
3.3 Formalization of KAOS Achieve Goals
If we refer to the concepts of guard and postcondition that exist in Event-B, a KAOS goal
can be considered as a postcondition of the system, since it means that a property must be
established. The crux of our formalization is to express each KAOS goal as a Event-B event,
where the action represents the achievement of the goal. Then, we will use the Event-B
refinement relation and additional custom-built proof obligations to derive all the subgoals
of the system by means of Event-B events. Let us consider now that the Achieve goal G is
refined into two sub-goals G1 and G2 as shown in Figure 1.
Figure 1: Example of KAOS goal model
Each level i (i ∈ [0..n]) is represented in the hierarchy of the KAOS goal graph as an EventB model Mi that refines the model Mi−1 related to the level i − 1. Moreover, we represent
A. Matoussi, F. Gervais, R. Laleau. An Event-B formalization of KAOS goal refinement patterns 6
MACHINE Abstract M0 MACHINE First M1
REFINES Abstract M0
EvG ∆= EvG1 ∆=
when G-Guard when G1-Guard
then G-PostCond then G1-PostCond
end end
EvG2 ∆=
when G2-Guard
then G2-PostCond
end
(a) Abstract Model M0 (b) Refinement Model M1
Figure 2: Overview of the Event-B representation of the KAOS goal model
each goal as a Event-B event where: (i) the current condition of this goal is considered as the
guard; (ii) the then part encapsulates the target condition of this goal (see Figure 2).
One may wonder how the temporal characteristic (”sooner or later” keywords) is expressed
in our Event-B formalization, as this should imply to specify the concept of time. However,
J.R Abrial [2] explains that the time dimension does not need to be explicitly considered
at the most abstract levels of a specification, and advocates to use events themselves to
express this dimension. Thus, in Event-B, each observable event is so small (atomic) that
its execution can be considered to take no time, and then only one event can take place
within one unit of time. Consequently, when the unit is large, the corresponding time is very
abstract, and when the unit is small, the corresponding time is very concrete. In other words,
time is stretched when moving from an abstraction to its refinement. This stretching reveals
some ”time details”, some new events. This interpretation suits well to the definition of an
Achieve goal. However, if deadlines were associated with such a goal, we would be obliged
to introduce new variables to model time. This point will be considered in future work when
non-functional goals will be formalized.
In the next sections, we propose an Event-B semantics for each KAOS refinement pattern
associated to Achieve goals. Based on the classical set of inference rules from Event-B [3],
we have identified the systematic proof obligations for each KAOS goal refinement pattern.
In that aim, the following outline will be used for describing each refinement pattern:
1. Description of the KAOS pattern: We give a short informal definition of the KAOS
goal refinement pattern.
2. Formal semantics of the pattern: We propose to give an Event-B semantics to the
KAOS goal refinement pattern as follows:
(a) Formal definition: We present the Event-B semantics given to each pattern by
constructing some set-theoretic mathematical models based on trace semantics of
Event-B developments [3].
(b) Proof obligations identification: We present some informal and formal arguments defining exactly what we have to prove for each KAOS goal refinement
pattern.
3. Synthesis: We summarize the Event-B formalization of each KAOS goal refinement
pattern.
A. Matoussi, F. Gervais, R. Laleau. An Event-B formalization of KAOS goal refinement patterns 7
4 Expressing the milestone-driven goal refinement pattern in
Event-B.
4.1 Description of the KAOS pattern
The milestone-driven goal refinement pattern [4] refines an Achieve goal by introducing intermediate milestone states G1, ..., Gn for reaching a state satisfying the target condition (denoted by G-PostCond) from a state satisfying the current condition (denoted by G-Guard)
as shown in Figure 3 (with just two sub-goals).
Figure 3: Milestone-driven goal refinement pattern
The first sub-goal G1 is an Achieve goal with the milestone condition as target condition;
it states that sooner or later the milestone condition (denoted by G1-PostCond) must hold if
the specific current condition G1-Guard (which can be larger than the current condition GGuard of the parent goal) holds in the current state. The second sub-goal is an Achieve goal
as well; it states that sooner or later the specific target condition G2-PostCond (which can
be larger than the target condition G-PostCond of the parent goal) must hold if the specific
milestone condition G2-Guard (derived from G1-PostCond) holds in the current state.
4.2 Formal semantics of the pattern
Since the satisfaction of all the KAOS sub-goals (according to a specific order) implies the
satisfaction of the parent goal, the abstract event EvG is refined by the sequence of all the
new events (EvG1, EvG2).
4.2.1 Formal definition
We propose a syntactic extension of the Event-B refinement proof rule in order to provide a
way to refine an abstract event by a sequence of new events as follows:
(EvG1 ; EvG2) Refines EvG
It is necessary to define the different proof obligations associated to this new refinement
semantics. Our idea is to try to characterize more accurately refinement in terms of trace
comparisons [3, 23] as illustrated in the following diagram.
A. Matoussi, F. Gervais, R. Laleau. An Event-B formalization of KAOS goal refinement patterns 8
S
∆= {v|I(v) }
T
∆= {w| ∃ v . (I(v) ∧ J(v, w)) }
EvG ∆= {v 7→ v
′
|I(v) ∧ G-Guard(v) ∧ G-PostCond(v, v′
) }
EvG1
∆= {w 7→ t| ∃ v . (I(v) ∧ J(v, w)) ∧ G1-Guard(w) ∧ G1-PostCond(w, t) }
EvG2
∆= {t 7→ w
′
| ∃ v
′
. (I(v
′
) ∧ J(v
′
, w′
)) ∧ G2-Guard(t) ∧ G2-PostCond(t, w′
) }
r
∆= {w 7→ v|I(v) ∧ J(v, w)) }
r
−1 ∆= {v 7→ w|I(v) ∧ J(v, w)) }
Figure 4: Set theoretic representation of the discrete models
The Event-B proof obligations must establish the well-known ”forward simulation” [23]
which is sufficient to guarantee refinement. Indeed, this condition ensures that any trace of
the refined model (EvG1 ; EvG2) must also be a trace of the abstraction EvG as follows:
r
−1
; EvG1; EvG2 ⊆ EvG; r
−1
where r a total binary relation from the concrete set T to the abstract set S. r formalizes
the gluing invariant between the refined state and the abstract one. In order to translate
this condition, it suffices to link S, T, EvG, EvG1, EvG2 and r with this new formulation.
Based on the set theoretic representation of the discrete models [3], each event is defined by
means of its guard and before-after predicate as stated in Figure 4.
To well express this refinement semantics, we shall extend the above set theoretic representation (see Figure 4) by expressing the sequence EvG1; EvG2 as follows:
EvG1; EvG2
∆= {w 7→ w
′
| ∃ v . I(v) ∧ J(v, w) ∧ ∃ t . G1-Guard(w) ∧ G1-PostCond(w, t)∧
G2-Guard(t) ∧ G2-PostCond(t, w′
) }
This definition expresses that we start in a state w where only EvG1 could be fired; i.e.
its guard holds. Once this event is executed (its post-condition became true), we reach immediately another intermediate state t in which the event EvG2 can be fired. The execution
of this last event allow us to reach the state w
′
. Based on this last definition, the two parts
of the forward simulation condition can be stated as follows:
A. Matoussi, F. Gervais, R. Laleau. An Event-B formalization of KAOS goal refinement patterns 9
r
−1
; EvG1; EvG2
∆= {v 7→ w
′
| ∃ w . I(v) ∧ J(v, w) ∧ ∃ t.G1-Guard(w)∧
G1-PostCond(w, t) ∧ G2-Guard(t) ∧ G2-PostCond(t, w′
) }
EvG; r
−1 ∆= {v 7→ w
′
| ∃ v
′
.I(v) ∧ G-Guard(v) ∧ G-PostCond(v, v′
) ∧ I(v
′
) ∧ J(v
′
, w′
) }
4.2.2 Proof obligations identification
Based on the above forward simulation condition, we are going to give systematic rules
defining exactly what we have to prove for this pattern. In fact, the translation of the
forward simulation condition, namely r
−1
; EvG1; EvG2 ⊆ EvG; r
−1
, comes down to prove
the following sequent1
:
I(v) ∧ J(v, w) ∧ G1-Guard(w) ∧ G1-PostCond(w, t) ∧ G2-Guard(t) ∧ G2-PostCond(t, w′
)
⊢
I(v) ∧ G-Guard(v) ∧ ∃ v
′
.G-PostCond(v, v′
) ∧ I(v
′
) ∧ J(v
′
, w′
)
Driven by the definition of the milestone refinement pattern, let us suppose that the target
condition of the goal G1 implies the current condition of the goal G2 as follows:
G1-PostCond ⇒ G2-Guard (PO1)
This first proof obligation allows us to simplify the last sequent as follows:
I(v) ∧ J(v, w) ∧ G1-Guard(w) ∧ G1-PostCond(w, t) ∧ G2-
PostCond(t, w′
)
⊢
G-Guard(v) ∧ ∃ v
′
.G-PostCond(v, v′
) ∧ I(v
′
) ∧ J(v
′
, w′
)
(I)
We know exactly now which sequent we have to prove. It is the only required condition
to verify in order to ensure the consistency of the KAOS milestone-driven goal refinement
pattern. This sequent (I) can be split into two sequents ((I.1) and (I.2)) as follows by applying
the inference rule AND-R2
:
I(v)
J(v, w)
G1-Guard(w)
G1-PostCond(w, t)
G2-PostCond(t, w′
)
⊢
G-Guard(v)
(I.1)
1The symbol ⊢ is named the turnstile. The part situated on the left hand part of the turnstile denotes a
finite set of predicates called the hypotheses. The part situated on the right hand side of the turnstile denotes
a predicate called the goal [3].
2
It allows us to simplify conjunctive predicates appearing in the goal of a sequent (see [3]).
A. Matoussi, F. Gervais, R. Laleau. An Event-B formalization of KAOS goal refinement patterns 10
Inspired by the definition of the milestone goal refinement pattern given by [4], we ensure
that proving the following proof obligation is sufficient to prove the sequent:
G1-Guard ⇒ G-Guard (PO2)
I(v)
J(v, w)
G1-Guard(w)
G1-PostCond(w, t))
G2-PostCond(t, w′
)
⊢
∃ v
′
.(G-PostCond(v, v′
) ∧ I(v
′
) ∧ J(v
′
, w′
))
(I.2)
As for the the sequent (I.1), we guarantee that proving the following proof obligation is
sufficient to prove the sequent:
G2-PostCond ⇒ G-PostCond (PO3)
4.3 Synthesis
The milestone-driven goal refinement pattern is expressed in Event-B by considering that the
abstract event EvG is refined by the sequence of all the new events ( EvG1 and EvG2).
For that, we have presented the different sufficient proof obligations (not necessary) associated to the syntactic extension of the refinement of Event-B that provides a way to refine
an abstract event by a sequence of new events as follows. Of course, these different proof
obligations must be in practice considered under an additional hypothesis containing a model
invariant I, a gluing invariant J and a collection P of axioms constraining constants and sets.
• The ordering constraint (PO1) expresses the ”milestone” characteristic between the
Event-B events.
• The guard strengthening (PO2) ensures that the concrete guard is stronger than the
abstract one. In other words, it is not possible to have the concrete version enabled
whereas the abstract one would not. The term “stronger” means that the concrete
guard implies the abstract guard.
• The correct refinement (PO3) ensures that the sequence of concrete events transforms
the concrete variables in a way which does not contradict the abstract event.
A. Matoussi, F. Gervais, R. Laleau. An Event-B formalization of KAOS goal refinement patterns 11
5 Expressing the AND goal refinement in Event-B
5.1 Description of the KAOS pattern
An Achieve goal G is AND refined into two (or more) sub-goals if the conjunction of the
sub-goals is sufficient to establish the satisfaction of the parent goal G as shown in Figure 5
(with just two sub-goals).
Figure 5: AND goal refinement
In Figure 5, the AND-refinement link expresses that the goal G is satisfied by satisfying
the goal G1 and the goal G2.
Restriction: In KAOS nothing is said on the execution order of the different sub-goals.
The formalization in Event-B reveals that this can constitute a serious problem when these
sub-goals manipulate shared variables. Let us explain more this problem by supposing that
a shared variable x is modified by two sub-goals (G1 and G2). G1 increases x by 3 and G2
divides x by 2. Hence, it is clear that the execution order of the two sub-goals is important.
To completely avoid this problem, a sufficient condition is to force a AND refinement to
manipulate only disjoint set of variables. This strong solution is quite close to other solutions
adopted by a lot of researchers such as J.R Abrial [1] with the parallel behavior concept. If
there is a AND refinement with shared variables, we propose to transform this form of AND
into a “milestone” refinement, and then to explicitly specify the order of modifications on
the shared variables.
5.2 Formal semantics of the pattern
As for the milestone-driven tactic, the satisfaction of all the KAOS sub-goals implies the
satisfaction of the parent goal. However, the execution of these new events must not necessary follows a specific order. Hence, our idea (inspired from Process Algebra [21]) is that
these events (EvG1 , EvG2) are executed in an arbitrary order: either EVG1;EVG2
or EVG2;EVG1. This corresponds to the semantics of the interleave operator in process
algebra.
5.2.1 Formal definition
We propose a syntactic extension of the Event-B refinement proof rule in order to refine an
abstract event by the interleaving of all the new events as follows:
(EvG1 ||| EvG2) Refines EvG
As for the milestone-driven tactic, we try to characterize more accurately refinement in
terms of trace comparisons as illustrated in the following diagram.
A. Matoussi, F. Gervais, R. Laleau. An Event-B formalization of KAOS goal refinement patterns 12
We have thus the following ”forward simulation” proof to perform in order to ensure that
any trace of a refined model must also be a trace of the abstraction:
r
−1
; (EvG1|||EvG2) ⊆ EvG; r
−1
To well express this new refinement semantics, we shall extend the above set theoretic
representation (see Figure 4) by expressing the interleaving EvG1|||EvG2 as follows:
EvG1|||EvG2
∆= {w 7→ w
′
| ∃ v . I(v) ∧ J(v, w) ∧ (G1-Guard(w) ∧ G2-Guard(w))∧
((G1-Guard(w) ∧ G2-Guard(w)) ⇒ G1-PostCond(w, w′
) ∧ G2-PostCond(w, w′
))}
This definition expresses that we start in a state w where both EvG1 and EvG2 can be
fired; i.e. their guards holds. Hence, we can achieved another state w
′
iff both events are
executed. Based on this last definition, the two parts of the forward simulation condition can
be stated as follows:
r
−1
; (EvG1|||EvG2) ∆= {v 7→ w
′
| ∃ w . I(v) ∧ J(v, w) ∧ (G1-Guard(w) ∧ G2-Guard(w))∧
((G1-Guard(w) ∧ G2-Guard(w)) ⇒ G1-PostCond(w, w′
) ∧ G2-PostCond(w, w′
))}
EvG; r
−1 ∆= {v 7→ w
′
| ∃ v
′
.I(v) ∧ G-Guard(v) ∧ G-PostCond(v, v′
) ∧ I(v
′
) ∧ J(v
′
, w′
) }
5.2.2 Proof obligations identification
Based on the above forward simulation condition, we are going to give systematic rules
defining exactly what we have to prove in order to ensure that a concrete event indeed
refines its abstraction. Indeed, the translation of the forward simulation condition, namely
r
−1
; (EvG1|||EvG2) ⊆ EvG; r
−1
, requires to prove the following sequent:
I(v) ∧ J(v, w) ∧ (G1-Guard(w) ∧ G2-Guard(w))∧
((G1-Guard(w) ∧ G2-Guard(w)) ⇒ G1-PostCond(w, w′
) ∧ G2-PostCond(w, w′
))
⊢
I(v) ∧ G-Guard(v) ∧ ∃ v
′
.G-PostCond(v, v′
) ∧ I(v
′
) ∧ J(v
′
, w′
)
This sequent must be simplified using inference rule IMP-L3 as follows:
I(v) ∧ J(v, w)∧
(G1-Guard(w) ∧ G2-Guard(w))∧
(G1-PostCond(w, w′
) ∧ G2-PostCond(w, w′
))
⊢
I(v) ∧ G-Guard(v) ∧ ∃ v
′
.G-PostCond(v, v′
) ∧ I(v
′
) ∧ J(v
′
, w′
)
3
It allows us to simplify implicative predicates appearing in the hypothesis part of a sequent (see [3]).
A. Matoussi, F. Gervais, R. Laleau. An Event-B formalization of KAOS goal refinement patterns 13
The definition of the AND goal refinement pattern given by [4] help us to detect the sufficient
proof obligations in order to prove the last sequent; i.e. we ensure that proving these proof
obligations is sufficient to prove the sequent:
G1-Guard ⇒ G-Guard (PO1)
G2-Guard ⇒ G-Guard (PO2)
(G1-PostCond ∧ G2-PostCond) ⇒ G-PostCond (PO3)
5.3 Synthesis
The AND goal refinement pattern is expressed in Event-B by considering that the abstract
event EvG is refined by the interleaving of all the new events ( EvG1 and EvG2). For that,
we have presented the sufficient proof obligations associated to the syntactic extension of the
refinement of Event-B that provides a way to refine an abstract event by an interleaving of
new events as follows:
• The guard strengthening (PO1, PO2) ensures that the concrete guard is stronger
than the abstract guard of EvG. The concrete guard of EvG1 ||| EvG2 can be either
G1-Guard (if we execute EvG1 at first) or G2-Guard (if we execute EvG2 at first).
• The correct refinement (PO3) ensures that the interleaving of concrete events EvG1
||| EvG2 transforms the concrete variables in a way which does not contradict the
abstract event.
Based on the description of all the possible behaviors, we give in Appendix A another
trace semantics that allows us to discover the above proof obligations. However, the new
semantics requires to prove the associativity condition if we process more than two events.
A. Matoussi, F. Gervais, R. Laleau. An Event-B formalization of KAOS goal refinement patterns 14
6 Expressing the OR goal refinement in Event-B
6.1 Description of the KAOS pattern
An Achieve goal G is OR refined into two sub-goals G1 and G2 if only either (not both) of
its sub-goals is achieved. A typical OR is shown in Figure 6 (with just two sub-goals).
Figure 6: OR goal refinement pattern
In each sub-goal, a sub-goal target condition must be reached in order to reach the target
condition G-PostCond of the parent goal. Notice that this goal refinement pattern refinement
introduces a certain kind of bounded non-determinism that will be resolved further in the
implementation phase.
6.2 Formal semantics of the pattern
6.2.1 Formal definition
Since the satisfaction of exactly one KAOS sub-goal implies the satisfaction of the parent
goal, we propose to refine the abstract event EvG as follows:
(EvG1 XOR EvG2) Refines EvG
This exclusive-OR refinement can be see as an inclusive OR refinement with an additional
exclusivity characteristic as follows:
(EvG1 XOR EvG2) Refines EvG =

(EvG1 OR EvG2) Refines EvG
Exclusivity characteristic (OR.def )
As for the other refinement patterns, the first part of the last definition (OR.def ) can be
characterized in terms of trace comparisons as illustrated in the following diagram.
We have thus the following ”forward simulation” proof to perform in order to ensure that
any trace of a refined model must also be a trace of the abstraction:
r
−1
; (EvG1 OR EvG2) ⊆ EvG; r
−1
A. Matoussi, F. Gervais, R. Laleau. An Event-B formalization of KAOS goal refinement patterns 15
To well express this new refinement semantics, we shall extend the above set theoretic
representation (see Figure 4) by expressing EvG1 OR EvG2 as follows:
EvG1 OR EvG2
∆= {w 7→ w
′
| ∃ v . I(v) ∧ J(v, w) ∧ (G1-Guard(w) ∧ G2-Guard(w))∧
((G1-Guard(w) ⇒ G1-PostCond(w, w′
)) ∨ (G2-Guard(w) ⇒ G2-PostCond(w, w′
)))}
This definition expresses that we start in a state w where both EvG1 and EvG2 can be
fired; i.e. their guards holds. Hence, we can achieved the state w
′
if at least one event is
executed. Notice that the ”exclusive” characteristic between these two events is an additional
constraint that will be considered further in the work. Based on this last definition, the two
parts of the forward simulation condition can be stated as follows:
r
−1
; (EvG1 OR EvG2) ∆= {v 7→ w
′
| ∃ w . I(v) ∧ J(v, w) ∧ (G1-Guard(w) ∧ G2-Guard(w))∧
((G1-Guard(w) ⇒ G1-PostCond(w, w′
)) ∨ (G2-Guard(w) ⇒ G2-PostCond(w, w′
)))}
EvG; r
−1 ∆= {v 7→ w
′
| ∃ v
′
.I(v) ∧ G-Guard(v) ∧ G-PostCond(v, v′
) ∧ I(v
′
) ∧ J(v
′
, w′
) }
6.2.2 Proof obligations identification
The translation of the forward simulation condition, namely r
−1
; (EvG1 OR EvG2) ⊆
EvG; r
−1
, requires to prove the following sequent:
I(v) ∧ J(v, w) ∧ (G1-Guard(w) ∧ G2-Guard(w))∧
((G1-Guard(w) ⇒ G1-PostCond(w, w′
))∨(G2-Guard(w) ⇒ G2-PostCond(w, w′
)))
⊢
I(v) ∧ G-Guard(v) ∧ ∃ v
′
.G-PostCond(v, v′
) ∧ I(v
′
) ∧ J(v
′
, w′
)
The application of the inference rules: OR-L 4
then IMP-L then MON 5 allows us to simplify
the above sequent by obtaining these two following sequents :
I(v) ∧ J(v, w) ∧ G1-Guard(w) ∧ G1-PostCond(w, w′
)
⊢
I(v) ∧ G-Guard(v) ∧ ∃ v
′
.G-PostCond(v, v′
) ∧ I(v
′
) ∧ J(v
′
, w′
)
The definition of the OR goal refinement pattern given by [4] help us to detect the sufficient
proof obligations in order to prove the last sequent; i.e. we ensure that proving these proof
obligations is sufficient to prove the sequent:
G1-Guard ⇒ G-Guard (PO1)
G1-PostCond ⇒ G-PostCond (PO2)
4
It corresponds to the classical technique of a proof by cases. More precisely, in order to prove a goal under
a disjunctive assumption P ∨ Q, it is sufficient to prove independently the same goal under assumption P and
also under assumption Q(see [3]).
5
It says that in order to have a proof of goal G under the two sets of assumptions H1 and H2, it is sufficient
to have a proof of G under H1 only (see [3]).
A. Matoussi, F. Gervais, R. Laleau. An Event-B formalization of KAOS goal refinement patterns 16
I(v) ∧ J(v, w) ∧ G2-Guard(w) ∧ G2-PostCond(w, w′
)
⊢
I(v) ∧ G-Guard(v) ∧ ∃ v
′
.G-PostCond(v, v′
) ∧ I(v
′
) ∧ J(v
′
, w′
)
Similarly, the second sequent can be proved by proving the following proof obligations:
G2-Guard ⇒ G-Guard (PO3)
G2-PostCond ⇒ G-PostCond (PO4)
The second part of the definition (OR.def ) that express the exclusive characteristic of the
OR refinement consists to prove these two proof obligations:
G1-PostCond ⇒ ¬G2-Guard (PO5)
G2-PostCond ⇒ ¬G1-Guard (PO6)
6.3 Synthesis
The OR goal refinement pattern is expressed in Event-B by considering that each concrete
event (EvG1 or EvG2) indeed refines its abstraction. However, we require an additional
proof obligations ensuring that it is the realization of exactly one sub-goal which allow the
realization of the parent goal:
• The guard strengthening (PO1, PO3) ensures that the concrete guard (G1-Guard or
G2-Guard) is stronger than the abstract guard of EvG.
• The correct refinement (PO2, PO4) ensures that each concrete event (EvG1 or
EvG2) transforms the concrete variables in a way which does not contradict the abstract event EvG.
• The ”exclusive” characteristic (PO5, PO6) ensures that only one event (either EvG1
or EvG2) but not both can be executed.
Most of these proof obligations could be discharged by the current version of the Rodin
automatic theorem prover [28]. In fact, this Event-B refinement semantics is quite close to
the same one proposed by Rodin if we consider that each event refines the abstract event
EvG.
A. Matoussi, F. Gervais, R. Laleau. An Event-B formalization of KAOS goal refinement patterns 17
7 Some other KAOS goal refinement patterns
The following patterns [4] introduce case conditions that guide the refinement of the parent
goal.
7.1 The decomposition-by-case pattern.
This pattern [4] is applicable to behavioral Achieve goals where different cases can be identified for reaching the target condition G-Post . The cases must be disjoint and cover the
entire state space. As shown in Figure 7, this refinement requires two domain properties
(color shape), one stating the disjointness and coverage property of the case condition Case1
and Case2, the other stating that the disjunction of the specific target conditions must imply
the target condition of the parent goal. The completeness of the AND-refinement is derivable
by use of those two domain properties.
Figure 7: Decomposition-by-cases pattern
The study of this pattern reveals that A. Van Lamsweerde [4] gives it a very large definition. For this reason, we restrict the condition of using this pattern. Hence, we recommend
requirements engineers to use this kind of pattern if we process just one element in the Achieve
goal. In this case, we must verify that the cases are disjoint and cover the entire state space.
This pattern can be seen as a special case of the OR refinement pattern since G1-Guard
and G2-Guard are very large than G-Guard and can encapsulate the different cases Case1
and Case2; i.e. these case are hidden in the guards G1-Guard and G2-Guard. Moreover,
the proof obligations associated to the OR refinement allows us to express the two domain
properties.
In the set-theoretic case, this pattern is not appropriate and have no sense. It is more
judicious to use the AND refinement pattern (with disjoint variable) that can express what
we want since we consider more large guards of sub-goals (G1-Guard and G2-Guard are very
large than G-Guard). In fact, these guards can encapsulate the different cases Case1 and
Case2.
Consequently, we don’t need to formalize the decomposition-by-case pattern in Event-B
since our basic patterns allow us to express such pattern.
7.2 The guard-introduction pattern.
This pattern is a case-driven refinement pattern applicable to behavioral Achieve goals where
a guard condition Condition must necessarily be set for reaching the target condition GPost. As shown in Figure 8, the first sub-goal of the Achieve goal states that the target
condition must be reached from a current condition G-Guard where in addition the guard
A. Matoussi, F. Gervais, R. Laleau. An Event-B formalization of KAOS goal refinement patterns 18
condition Condition must hold. The second sub-goal states that the guard condition Condition must be reached as a target from that current condition G-Guard. The third subgoal states that the current condition G-Guard must always remain true unless the target
condition G-Post of the parent goal is reached.
Figure 8: Guard-introduction pattern
If we switch the first two sub-goals, we remark that we obtain easily the milestone refinement pattern; i.e. if we consider of course that G-Guard must always remain true unless
G-Post is reached (the third sub-goal). Hence, the guard-introduction pattern is a special
case of the milestone refinement pattern. So, we don’t need to formalize this pattern in
Event-B.
8 Case Study
For the case study, we consider the specification of a localization software component which is
a critical part of a land transportation system. Many positioning systems have been proposed
over the last years. GPS, one of the most widely used positioning system, is perhaps the
best-known. This system belongs to the GNSS (Global Navigation Satellite Systems) family
which also regroups GALILEO or GLONASS. Positioning systems are often dedicated to a
particular environment; the GNSS technology, for example, generally does not work indoors.
To resolve these problems, numerous alternatives relying on very different technologies have
arisen. Those last years, Wireless LAN such as IEEE 802.11 networks have been considered by
numerous location systems. These systems all use the radio signal strength to determine the
physical location. Localization systems can therefore be designed using various technologies
like wireless personal networks such as Wifi or Bluetooth [25, 26], sensors [27], GNSS repeaters
or visual landmarks.
The main difficulties when we develop a localization component is to find the correct
algorithm that merges positioning data and to take into account all the properties we have
to deal with. At this stage, we think that a semi-formal model will be very useful to have
guidelines on how to do. However, elaborating this semi-formal model is not necessarily an
easy task. Often, requirements engineers need to search through preliminary documents in
order to extract goals (key properties) using a number of heuristics (asking HOW and WHY
questions...) detailed in [4]. Figure 9 show the obtained KAOS goal model of a localization
component thanks to heuristics. For example, a HOW question about the goal G would
then lead to the goals G1, G2 and G3. This obtained KAOS goal model contains: (i) highlevel goals; (ii) refinement links denoted by a bubble linking the parent goal with an arrow,
and the child goals with regular lines; (ii) requirements and their software responsibility
A. Matoussi, F. Gervais, R. Laleau. An Event-B formalization of KAOS goal refinement patterns 19
agents; (iv) expectations and their environment responsibility agents (GPS, WIFI, sensor
and accelerometer).
Figure 9: KAOS goal model of a localization component
Let us start by the high level goal G which is defined as follows:
Goal G: Achieve [LocalizeVehicle]
InformalDef: The Cycab/vehicle must be localized.
We associate an Event-B model Localization, in Figure 10, to this most abstract level of the
hierarchy of the KAOS goal graph. In this Event-B model, we will have an event called LocalizeVehicle that will translate the goal G; i.e. it describes the ”property” of the goal G, in
terms of generalized substitutions. Localizing a vehicle consists in obtaining an estimated loc
which is a couple of latitude and longitude. At this level of abstraction, it is not necessary to
precise the way this information is obtained. Thus we use the non-deterministic generalized
substitution through the symbol :∈ . It specifies an unbounded choice and estimated loc can
take any value in the sets (LATITUDE \ {null}) and (LONGITUDE \ {null}). The null
value serves just to initialize the system through the initialization event6
. Notice that at
this abstraction level, the event LocalizeVehicle can always occurs. Hence, its guard is set
to true. Sets in uppercase are abstract sets used to type the variables. They are described
in the Event-B context TypeSets (see in Figure 11).
6The initialization part, variables, invariants and contexts are manually completed by the designer.
A. Matoussi, F. Gervais, R. Laleau. An Event-B formalization of KAOS goal refinement patterns 20
MACHINE Localization
SEES TypeSets
VARIABLES
estimated loc
INVARIANTS
inv1 : estimated loc ∈ LATITUDE × LONGITUDE
EVENTS
Initialisation
begin
act1 : estimated loc := null 7→ null
end
Event LocalizeVehicle =b
begin
act1 : estimated loc :∈ (LATITUDE \ {null}) × (LONGITUDE \ {null})
end
END
Figure 10: The abstract model
CONTEXT TypeSets
SETS
SUBCOMPONENTS
SUBSENSORS
CONSTANTS
gps, wifi, LATITUDE, LONGITUDE
null, speed, accel
AXIOMS
axm1 : partition(SUBCOMPONENTS, {gps}, {wifi})
axm2 : LATITUDE = N ∪ {null}
axm3 : LONGITUDE = N ∪ {null}
axm4 : partition(SUBSENSORS, {speed}, {accel})
END
Figure 11: The Event-B context TypeSets
A. Matoussi, F. Gervais, R. Laleau. An Event-B formalization of KAOS goal refinement patterns 21
8.1 First refinement
The goal G is refined into three sub-goals according to the milestone goal refinement pattern:
Goal G1: Achieve [CaptureRawLocalizations]
InformalDef: Firstly, several sets of raw localization data are captured by using different technologies.
Goal G2: Achieve [ValidateData]
InformalDef: Then, all the sets of raw localization data will be validated and controlled.
Goal G3: Achieve [MergeData]
InformalDef: Finally, all the validated data will be merged in order to obtain the
final localization.
Similarly, we associate an Event-B refinement model Localization1, in Figure 12, to this
first level of the hierarchy of the KAOS goal graph. The sub-goals G1, G2 and G3 are represented by three Event-B events CaptureRawLocalizations, ValidateData and MergeData, respectively. The first one returns a set of couples (latitude, longitude), one for each
component used for localizing a vehicle. The second one validates the returned set of couples
by choosing the acceptable values. The final one returns the final localization calculated from
the returned values of the event MergeData.
In Event-B, often the information about such event ordering has to be embedded into
guards and event actions with the downside of extra model variables. As we said previously,
we have chosen to explicitly reproduce KAOS goal ordering in an Event-B model by proposing
a syntactic extension of the Event-B refinement proof rule in order to provide a way to refine
an abstract event by a sequence of new events. Hence, the abstract event LocalizeVehicle
is refined as follows:
(CaptureRawLocalizations; ValidateData; MergeData) Refines LocalizeVehicle
In addition to the feasibility proof obligation7
, this kind of refinement requires to discharge
these different proof obligations:
• Two ordering constraints express the ”milestone” characteristic between the Event-B
events. These two proof obligation are discharged: (i) the action of CaptureRawLocalizations implies the guard of ValidateData (ii) the action of ValidateData implies the guard of MergeData.
• One “guard strengthening” is also discharged since the first event (CaptureRawLocalizations)
in the sequence has a guard (true) that implies the abstract guard (true).
• One “correct refinement” is also proved since the last event (MergeData) in the sequence has a postcondition that implies the abstract postcondition under the gluing
invariant inv4. Hence, the sequence of concrete events transforms the concrete variables in a way which does not contradict the abstract event.
7
It ensures that each event must also be feasible, in a sense that an appropriate new state v
′ must exist for
some given current state v.
A. Matoussi, F. Gervais, R. Laleau. An Event-B formalization of KAOS goal refinement patterns 22
MACHINE Localization1
REFINES Localization
SEES TypeSets
VARIABLES
estimated loc, subcomponents loc, validated loc, merged loc
INVARIANTS
inv1 : subcomponents loc ∈ SUBCOMPONENTS → (LATITUDE × LONGITUDE)
inv2 : validated loc ∈ SUBCOMPONENTS 7→ (LATITUDE × LONGITUDE)
inv3 : merged loc ∈ LATITUDE × LONGITUDE
inv4 : estimated loc = merged location
EVENTS
Initialisation
begin
act1 : estimated loc := null 7→ null
act4 : subcomponents loc :∈ SUBCOMPONENTS → ({null} × {null})
act3 : validated loc :∈ SUBCOMPONENTS → ({null} × {null})
act5 : merged loc := null 7→ null
end
Event CaptureRawLocalizations =b
begin
act1 : subcomponents loc :∈ SUBCOMPONENTS →((LATITUDE \ {null})×(LONGITUDE \
{null}))
end
Event ValidateData =b
when
grd1 : subcomponents loc ∈ SUBCOMPONENTS → ((LATITUDE \ {null}) × (LONGITUDE \
{null}))
then
act1 : validated loc :∈ P1 (subcomponents loc)
end
Event MergeData =b
when
grd1 : validated loc ∈ P1 (subcomponents loc)
grd2 : subcomponents loc ∈ SUBCOMPONENTS → ((LATITUDE \ {null}) × (LONGITUDE \
{null}))
then
act1 : merged loc :∈ (LATITUDE \ {null}) × (LONGITUDE \ {null})
end
END
Figure 12: First Event-B refinement model
A. Matoussi, F. Gervais, R. Laleau. An Event-B formalization of KAOS goal refinement patterns 23
8.2 Second refinement
Now, we consider the second level of the hierarchy of the KAOS goal graph. In the same
way, a refinement Event-B model Localization2, in Figure 13, is created and must refine the
previous model Localization1. This second refinement Event-B model will encapsulate two
KAOS refinement patterns:
8.2.1 Second refinement: Applying the AND goal refinement pattern
The goal G1 is AND-refined into two sub-goals. This refinement specifies the kind of technology used to obtain localization data.
Goal G1.1: Achieve [UseGPS]
InformalDef: A GPS system is used.
Goal G1.2: Achieve [UseWIFI]
InformalDef: A wireless technique is used.
The sub-goals G1.1 and G1.2 are represented by two Event-B events UseGPS and UseWIFI
by using the same transcription rules as for the event LocalizeVehicle in the abstract model.
Since the goal G1 is refined into two sub-goals G1.1 and G1.2 according to the AND goal
refinement pattern, the execution of the corresponding new events must not necessary follows a specific order. As we said previously, our idea (inspired from Process Algebra [21])
is that these events (UseGPS and UseWIFI are executed in an arbitrary order: either
UseGPS;UseWIFI or UseWIFI;UseGPS as follows (of course, we must ensure that this
AND refinement manipulate only disjoint set of variables):
(UseGPS ||| UseWIFI) Refines CaptureRawLocalizations
In addition to the feasibility proof obligation, the following proof obligations must be
discharged in order to prove such refinement:
• Two “guard strengthening” are discharged since the concrete guard of (UseGPS |||
UseWIFI) implies the abstract guard (true) of CaptureRawLocalizations. In fact,
the concrete guard is always true (if we execute UseGPS at first or if we execute
UseWIFI at first).
• One “correct refinement” is also proved since the conjunction of the two concrete postconditions implies the abstract postcondition under the gluing invariant inv3. Hence,
this ensures that subcomponents loc is a total function (see the postcondition of the
abstract event CaptureRawLocalizations).
8.2.2 Second refinement: Applying the milestone goal refinement pattern
On the other hand, the goal G2 is refined into two sub-goals according to the milestone-driven
tactics:
Goal 2.1: Achieve [CaptureRelativeLocalizations]
InformalDef: At first, several sets of relative localization data are captured by using
A. Matoussi, F. Gervais, R. Laleau. An Event-B formalization of KAOS goal refinement patterns 24
MACHINE Localization2
REFINES Localization1
SEES TypeSets
VARIABLES
estimated loc, subcomponents loc, validated loc, merged loc
gps loc, wifi loc, sensors loc, kept loc
INVARIANTS
inv1 : gps loc ∈ {gps} → (LATITUDE × LONGITUDE)
inv2 : wifi loc ∈ {wifi} → (LATITUDE × LONGITUDE)
inv3 : subcomponents loc = gps loc ∪ wifi loc
inv4 : sensors loc ∈ SUBSENSORS 7→ (LATITUDE × LONGITUDE)
inv5 : kept loc ∈ SUBCOMPONENTS 7→ (LATITUDE × LONGITUDE)
inv6 : validated loc = kept loc
EVENTS
Initialisation
begin
act1 : estimated loc := null 7→ null
act4 : subcomponents loc :∈ SUBCOMPONENTS → ({null} × {null})
act3 : validated loc :∈ SUBCOMPONENTS → ({null} × {null})
act5 : merged loc := null 7→ null
act7 : gps loc :∈ {gps} → ({null} × {null})
act6 : wifi loc :∈ {wifi} → ({null} × {null})
act8 : sensors loc :∈ SUBSENSORS → ({null} × {null})
act9 : kept loc :∈ SUBCOMPONENTS → ({null} × {null})
end
Event UseGPS =b
begin
act1 : gps loc :∈ {gps} → ((LATITUDE \ {null}) × (LONGITUDE \ {null}))
end
Event UseWIFI =b
begin
act1 : wifi loc :∈ {wifi} → ((LATITUDE \ {null}) × (LONGITUDE \ {null}))
end
Event CaptureRelativeLocalizations =b
when
grd1 : subcomponents loc ∈ SUBCOMPONENTS → ((LATITUDE \ {null}) × (LONGITUDE \
{null}))
then
act1 : sensors loc : |(sensors loc′ ∈ SUBSENSORS 7→ ((LATITUDE \ {null}) ×
(LONGITUDE \ {null}))) ∧ sensors loc′
6= ∅
end
Event FilterData =b
when
grd1 : sensors loc ∈ SUBSENSORS 7→ ((LATITUDE \ {null}) × (LONGITUDE \ {null})) ∧
sensors loc 6= ∅
then
act1 : kept loc :∈ P1 (subcomponents loc)
end
END
Figure 13: Second Event-B refinement model
A. Matoussi, F. Gervais, R. Laleau. An Event-B formalization of KAOS goal refinement patterns 25
different technologies.
Goal 2.2: Achieve [FilterData]
InformalDef: Then, all the sets of raw localization data will be filtered.
All these subgoals are translated into new events using the same rules as for LocalizeVehicle in the abstract Event-B model. As for the first refinement, the abstract event ValidateData is refined by the sequence of all the new events (CaptureRelativeLocalizations,
FilterData) as follows:
(CaptureRelativeLocalizations ; FilterData) Refines ValidateData
We have also discharged the different proof obligations related to the milestone refinement
such as the ordering constraint, the “guard strengthening” and the “correct refinement”.
8.3 Third refinement
The goal G2.1 is OR-refined in two subgoals:
Goal G2.1.1: Achieve [UseSpeedSensor]
InformalDef: The Cycab may use a speed sensor system.
Goal G2.1.2: Achieve [UseAccelerometer]
InformalDef: Or, it may use the accelerometer system.
Similarly, a refinement Event-B model Localization3, in Figure 14, is associated to this
third level of the hierarchy of the KAOS goal graph. All these subgoals are transformed
into new Event-B events (UseSpeedSensor, UseAccelerometer) by using the same transcription rules as for the event LocalizeVehicle in the abstract Event-B model. Since the
satisfaction of exactly one KAOS sub-goal implies the satisfaction of the parent goal, we
propose to refine the abstract event CaptureRelativeLocalizations as follows:
(UseSpeedSensor XOR UseAccelerometer) Refines CaptureRelativeLocalizations
As we said previously, this Event-B refinement semantics is quite close to the same one
proposed by Rodin [28] if we consider that each event refines the abstract event. Hence,
the proof obligations (“guard strengthening” and “correct refinement”) could be discharged
by the current version of the Rodin automatic theorem prover [28]. However, we require
to express two additional proof obligations ensuring that only one event (either UseSpeedSensor or UseAccelerometer) but not both can be executed. These two proof obligations
are discharged since (i) the postcondition of UseSpeedSensor forbids the guard of UseAccelerometer to be triggered; (ii) the postcondition of UseAccelerometer forbids the
guard of UseSpeedSensor to be triggered.
This is the last step of refinement since all the goals are either requirements or expectations.
This allows to obtain the abstract Event-B specification from the KAOS goal hierarchy. The
Event-B model is then further refined towards an implementation. It concerns only the
Event-B events corresponding to requirements assigned to a software argent. Otherwise,
an expectation (assigned to an external agent) is a property required on the environment
and it corresponding Event-B event will not be implemented in the software-to-be. More
precisely, in our case study, the Event-B events UseGPS, UseWIFI, UseSpeedSensor
A. Matoussi, F. Gervais, R. Laleau. An Event-B formalization of KAOS goal refinement patterns 26
MACHINE Localization3
REFINES Localization2
SEES TypeSets
VARIABLES
sensors loc, speed loc, accel loc
INVARIANTS
inv1 : speed loc ∈ {speed} → (LATITUDE × LONGITUDE)
inv2 : accel loc ∈ {accel} → (LATITUDE × LONGITUDE)
inv3 : (sensors loc = speed loc) ∨ (sensors loc = accel loc) ∨ (sensors loc = speedloc ∪ accel loc)
EVENTS
Initialisation
begin
act10 : speed loc :∈ {speed} → ({null} × {null})
act11 : accel loc :∈ {accel} → ({null} × {null})
end
Event UseSpeedSensor =b
refines CaptureRelativeLocalizations
when
grd1 : subcomponents loc ∈ SUBCOMPONENTS → ((LATITUDE \ {null}) × (LONGITUDE \
{null}))
grd2 : accel loc ∈ {accel} → ({null} × {null})
then
act1 : speed loc, sensors loc : |(speed loc′ ∈ {speed}→((LATITUDE\{null})×(LONGITUDE\
{null}))) ∧ sensors loc′ = speed loc′
end
Event UseAccelerometer =b
refines CaptureRelativeLocalizations
when
grd1 : subcomponents loc ∈ SUBCOMPONENTS → ((LATITUDE \ {null}) × (LONGITUDE \
{null}))
grd2 : speed loc ∈ {speed} → ({null} × {null})
then
act1 : accel loc, sensors loc : |(accel loc′ ∈ {accel}→((LATITUDE \ {null})×(LONGITUDE \
{null}))) ∧ sensors loc′ = accel loc′
end
END
Figure 14: Third Event-B refinement model
A. Matoussi, F. Gervais, R. Laleau. An Event-B formalization of KAOS goal refinement patterns 27
lat, long ← OPMerge =
VAR lat gps, long gps, lat wifi, long wifi, ponderation gps, ponderation wifi
IN
lat gps := get lat(gps loc)||
long gps := get long(gps loc) ||
ponderation gps := get pond(gps loc) ||
lat wifi := get lat(wifi loc)||
long wifi := get long(wifi loc) ||
ponderation wifi := get pond(wifi loc) ;
lat := ((lat gps ∗ ponderation gps) + (lat wifi ∗ ponderation wifi)) /
(ponderation gps + ponderation wifi) ||
long := ((long gps ∗ ponderation gps) + (long wifi ∗ ponderation wifi)) /
(ponderation gps + ponderation wifi)
END
Figure 15: The B operation OPMerge
and UseAccelerometer are not refined since they correspond to goals of type expectation.
They are implemented by hardware components (GPS, WIFI components...) in a vehicle.
Only the Event-B events FilterData and MergeData are further refined. For example,
the refinement of the Event-B event MergeData leads to a software that implements the
algorithm chosen to realize the fusion (thanks to the B operation OPMerge) as shown in
Figure 15. At first, this operation recovers all the raw localization data from both GPS and
WIFI thanks to the call of operations get lat and get long. Then, the call of the operation
get pond serves to verify if the returned values of GPS and WIFI are validated or no. If these
values are validated, so get pond returns a weighting value set to 1 (O otherwise). Finally, the
operation calculates the final latitude and longitude based on the different weighting values.
An interesting result is that the link between the B operation OPMerge and the abstract
Event-B event MergeData (see Figure 12) is ensured. While the abstract event MergeData
describes the properties that the final program must fulfill, the B operation OPMerge
describes the algorithm contained in the program. Hence, MergeData describes the way by
which we can eventually judge that the final program OPMerge is correct: (i) the call of
the operations get lat and get long ensures the second guard of MergeData; (ii) the call
of the operation get pond ensures the first guard of MergeData; (iii) the final result of the
operation OPMerge (lat, long) satisfies the post-condition of MergeData.
9 Discussion
One may wonder whether the formalization of KAOS target conditions as Event-B postconditions is adequate, since the execution of Event-B events is not mandatory. In accordance
with the Event-B semantics, all events whose guard is true can be performed and there is
necessarily one that will be performed. The choice between all these permitted events is made
non-deterministically.
Another important point is that up to now we have not studied the combination and the
interaction between the different KAOS refinement pattern. Consequently, we are not certain
to maintain all the proof obligations. For instance, the study of the global behavior of the
KAOS goal graph, thanks to trace semantics, reveals the following main problem caused by
the strong proof obligation PO1 of the milestone refinement pattern (see Section 4). Let us
explain this problem by considering the following goal model (Figure 16):
A. Matoussi, F. Gervais, R. Laleau. An Event-B formalization of KAOS goal refinement patterns 28
Figure 16: KAOS goal model
According to our Event-B refinement semantics, the events of the first refinement are
interleaved as follows:
ev1|||ev2|||ev3
Now, one of possible valid execution trace in the second refinement is :
ev1.1 ; ev2.1 ; ev2.2 ; ev1.2 ; ev3.1 ; ev3.2
The main problem is that the execution of ev2.1 or ev2.2 can imply that the guard of
ev1.2 become false. So, the proof obligation ”ordering constraint” of the milestone refinement
pattern (ev1.1-Post ⇒ ev1.2-Guard) can never be discharged in this case. Indeed, it is more
logical now to relax this proof obligation by updating it as follows:
✷(ev1.1-Post ⇒ ✸ ev1.2-Guard)
This property means that it is always the case that once ev1.1 − Post holds then ev1.2 −
Guard holds eventually. Notice that ev1.1−Post is only required to hold at the initialization
of this process. So, ev1.1−Post is not required to hold any more before achieved g1.2−Guard.
This semantic corresponds to the classical temporal operator [22] leadsto (❀) as follows:
ev1.1 − Post ❀ ev1.2 − Guard
Proof techniques as in the Event-B method cannot be used to verify such ordering constraints. It seems that model checking may be a good alternative to prove such properties
that contain a temporal operator. So, our aims is to mix “local” proof obligations (generated by the prover) with the proof generated by the model checker. We think that a tool like
ProB [15] may be a very useful tool since it helps to discover the faults, where the consistency
is not completely achieved by the B prover.
The interest of the proposed approach is that we can prove some properties of consistency
on the goal model by discharging the proof obligations generated by the Event-B refinement
process. Indeed, if we can discharge for instance the proof obligations of refinement of an
event EvG by either EvG1 or EvG2 it means that the OR decomposition of goal G is
correct. Otherwise, it means either that one or more expressions of the events (EvG, EvG1,
A. Matoussi, F. Gervais, R. Laleau. An Event-B formalization of KAOS goal refinement patterns 29
EvG2) are not correct or that sub-goals are missing or that the goal refinement pattern is
false. However, we can never ensure that the expression of the Event-B events corresponds
exactly to the expression of the related goal since this later is informal. This means that
proving the Event-B formal counterpart of the KAOS goal model is not sufficient in order
to validate the conformance of the specification to original requirements. For that, we can
use an animation technique to validate the derived formal specification and consequently its
semi-formal counterpart goal model against original customers’ requirements. This animation
step not only indicates deviations from original requirements right on the spot but also helps
fixing some specification errors. The reader can refer to [17] for more details. For that,
ProB [15] may be also a very useful validation tool since its automated animation facilities
allow users to animate their specifications.
10 Related work
Our proposed approach aims at expressing KAOS models with Event-B language by staying
at the same abstraction level. In the sequel, we outline a number of approaches that have
tried to bridge the gap between KAOS requirements model and formal methods. With the
best of our knowledge, they are the only work that deal with such a problematic.
The work of [13] presents a goal-oriented approach to elaborate a pertinent model and turn
it into a high quality abstract B machines. The scheme used by the authors for transforming
the KAOS requirements model to B is as follows: As agents are the active entities able to
perform operations, a B machine is associated with each KAOS agent. The agent attributes
and the operations arguments are represented by the sets, variables and constraints. All
maintain goals under the agent responsibility are translated as invariants of the corresponding
B machine. All the KAOS operations that an agent has to perform are represented by B
operations.
The authors of [9] provides means for transforming the security requirements model built
with KAOS to an equivalent one in B. This abstract B model is then refined using non-trivial
B refinements that generate design specifications conforming to the initial set of security
requirements. The authors consider each operational goal and each KAOS operation as a B
operation. Also, they consider that each KAOS object, related to the operational goals, is
B machine. So, The relationship among objects is captured using the B machine imports,
includes, uses, and sees clauses that allow one B machine to relate to or compose other B machines. KAOS domain properties are transformed to B invariants or pre-conditions related to
the corresponding B operations. The authors introduce the concept of goal achievement which
is reflected through the return values of the B operations that model KAOS goals. Hence,
each B operation corresponding to an operational goal returns a flag indicating whether the
goal implemented in this operation has been achieved or not.
We can also point out a work [12] proposing an automatically generator that transforms
an extend KAOS model into VDM++ specifications. The generator connects operations in
KAOS to those in VDM++, and entities in KAOS to objects or types in VDM++. The
generated specification contains implicit operations consisting of pre- and post-conditions,
inputs, and outputs of operations. However, this generated specifications require software
developers to add the body of operations in order to create explicit specifications.
The GOPCSD (Goal-oriented Process Control System Design) tool [8] is an adaptation
of the KAOS method that serves to analyze the KAOS requirements and generate B formal
specifications. The tool is used to construct the application requirements in the form of goalmodels by interacting with the user and importing library templates. Then, the requirements
A. Matoussi, F. Gervais, R. Laleau. An Event-B formalization of KAOS goal refinement patterns 30
are checked to enable the system engineer to debug and correct them. Finally, the requirements will be translated to B specifications. The generated specifications can be refined and
translated to executable code by a software engineer.
Recently, [16] presents a constructive verification-based approach that consists in linking
requirements, expressed as linear temporal logic formulae, to a system specification expressed
as an Event-B machine extended with the notion of obligations [14]. The source requirements
are included as verification assertions that can be model-checked by tools like ProB [15],
showing that the proposed specification indeed meets the system requirements.
Nevertheless, the reconciliation presented by all of these works remains partial because
they don’t consider all the parts of the KAOS goal model but only the requirements (operational goals). Consequently, the formal model does not include any information about the
non-operational goals and, more important, the type of goal refinement. In this paper, we
have explored how to cope with this problem using an approach that expresses the whole
KAOS goal model with a formal method like Event-B by staying at the same abstraction
level. Our approach can be considered as complementary to existing ones. Furthermore,
what we present can be very useful in practice to (i) systematically verify that all KAOS
requirements are represented in the Event-B model; (ii) systematically verify that each element in the Event-B model has a purpose in KAOS. Moreover, employing formal methods in
requirements engineering level allows us to detect anomalies when we use the goal refinement
pattern in a chaotic manner. Hence, formal methods offers a recommendation support to
requirement engineers in order to choose the appropriate goal refinement patterns.
11 Conclusion and further work
In this report, we have presented a constructive approach driven by goals showing that it
is possible to express KAOS goal models with Event-B. Even if formal methods are harder
to use and less widely applicable, we have shown that extending KAOS with more formality
provide much higher precision and richer forms of analysis. The main contribution of our
approach is that it establishes the first brick toward the construction of the bridge between
the non-formal and the formal worlds as narrow and concise as possible. This brick balances
the trade-off between complexity of rigid formality (Event-B method) and expressiveness of
semi-formal approaches (KAOS). However, a number of future research steps are ongoing.
Regarding the different KAOS goal model concepts, we need now to consider the translation
of the concepts of domain properties and non functional goals. We plan also to apply the
approach on a number of case studies. Moreover, it would be interesting to establish the
correspondence between the obtained Event-B representation of KAOS goal models and the
later phases of development. At tool level, we plan to develop a connector between a KAOS
toolkit and the RODIN open platform.
References
[1] J.R. Abrial. The B-Book: Assigning programs to meanings. Cambridge University Press,
1996.
[2] J.R. Abrial. Extending B without changing it (for developing distributed systems). The
First Conference on the B-Method,pages 169–190, IRIN, Nantes, France, 1996.
A. Matoussi, F. Gervais, R. Laleau. An Event-B formalization of KAOS goal refinement patterns 31
[3] J.R. Abrial. Modeling in Event-B: System and Software Engineering. Cambridge University Press, 2010.
[4] A. van Lamsweerde. Requirements Engineering: From System Goals to UML Models to
Software Specifications. Wiley, 2009.
[5] R. Darimont and A. van Lamsweerde. Formal Refinement Patterns for Goal-Driven
Requirements Elaboration. In SIGSOFT ’96, pages 179–190, San Francisco, California,
USA, October 1996. ACM SIGSOFT.
[6] P. Behm, P. Benoit, A. Faivre, and J.-M. Meynadier. METEOR : A successful application
of B in a large project. In FM ’99: Proceedings of the Wold Congress on Formal Methods
in the Development of Computing Systems, Volume I pages 369–387, 1999. Springer.
[7] F. Badeau and A. Amelot. Using B as a high level programming language in an industrial
project: Roissy val. In Proceedings of the 4th International Conference of B and Z Users
(ZB’05), pages 334–354, Guildford, UK, 2005. Springer.
[8] I. El-Madah and T. Maibaum. Goal-Oriented Requirements Analysis for Process Control
Systems Design. In MEMOCODE 2003, pages 45–46, Mont Saint-Michel, France, June
2003. IEEE Computer Society.
[9] R. Hassan and S. Bohner and S. El-Kassas and M. Eltoweissy. Goal-Oriented, B-Based
Formal Derivation of Security Design Specifications from Security Requirements. In
ARES 2008, pages 1443–1450, Barcelona, Spain, March 2008. IEEE Computer Society.
[10] A. Matoussi and F. Gervais and R. Laleau A First Attempt to Express KAOS Refinement
Patterns with Event B. In ABZ 2008, pages 338, London, UK, September 2008. Springer.
[11] A. Matoussi and F. Gervais and R. Laleau De KAOS vers Event-B : Approche dirig´ee
par les buts. In Actes de la conf´erence AFADL’2009, pages 71–86, Toulouse, France,
Janvier 2009.
[12] H. Nakagawa and K. Taguchi and S. Honiden. Formal Specification Generator for KAOS:
Model Transformation Approach to Generate Formal Specifications from KAOS Requirements Models. In ASE 2007, pages 531–532, Atlanta, Georgia, USA, November 2007.
ACM.
[13] C. Ponsard and E. Dieul. From Requirements Models to Formal Specifications in B. In
REMO2V’2006, Luxembourg, June 2006.
[14] J. Bicarregui and A. Arenas and B. Aziz and P. Massonet and C. Ponsard. Towards
Modelling Obligations in Event-B. In ABZ 2008, pages 181–194, London, UK, September
2008. Springer.
[15] M. Leuschel and M.J. Butler. ProB: A Model Checker for B. In K. Araki, S. Gnesi,
D. Mandrioli (eds), FME 2003: Formal Methods, LNCS 2805, pages 855–874, 2003.
Springer.
[16] B. Aziz and A. Arenas and J. Bicarregui and C. Ponsard and P. Massonet. From GoalOriented Requirements to Event-B Specifications. In In: First Nasa Formal Method
Symposium (NFM 2009) , Moffett Field, California , USA, April 2009.
A. Matoussi, F. Gervais, R. Laleau. An Event-B formalization of KAOS goal refinement patterns 32
[17] A. Mashkoor and A. Matoussi Towards Validation of Requirements Models. In 2nd
International Conference on Abstract State Machines (ASM), Alloy, B and Z (ABZ’10),
Orford, Canada, 2010. To appear.
[18] A. Mammar and R. Laleau A formal approach based on UML and B for the specification
and development of database applications. In Automated Software Engineering 13(4),
pages 497-528, 2006.
[19] C. Snook and M. Butler. UML-B: formal modelling and design aided by UML. In ACM
Transactions on Software Engineering and Methodology, 15(1), pp. 92-122, 2006.
[20] E. Yu. Towards Modeling and Reasoning Support for Early-Phase Requirements Engineering. In RE’97, pages 226-235, 1997. IEEE Computer Society.
[21] D. Sangiorgi. Locality and interleaving semantics in calculi for mobile processes. In
Theor. Comput. Sci., Volume 155, pages 39–83, 1996. Elsevier Science Publishers Ltd.
[22] Z. Manna and A. Pnueli. Adequate proof principles for invariance and liveness properties
of concurrent programs. Sci. Comput. Program., Volume 4, number 3, pages 257–289,
1984. Elsevier North-Holland, Inc.
[23] N. Lynch and F. Vaandrager. Forward and Backward Simulations - Part I: Untimed
Systems. Information and Computation Journal, Volume 121, pages 214–233, 1994.
[24] W.N. Robinson and S. Pawlowski. Surfacing Root Requirements Interactions from Inquiry Cycle Requirements Documents. In The Third IEEE International Conference on
Requirements Engineering (ICRE’98), pages 82–89, Colorado Springs, CO, USA, 1998.
IEEE Computer Society Press.
[25] J. Hallberg and M. Nilsson and K. Synnes. Positioning with bluetooth. In 10th Int.
Conference on Telecommunications (ICT’2003), pages 954—958, 2003.
[26] J.A. Royo and E. Mena and L.C. Gallego. Locating Users to Develop Location-Based
Services in Wireless Local Area Networks. In UCAmI2005, pages 471-478, Granada,
Spain, 2005.
[27] R.J. Orr and G.D. Abowd. The smart floor: A mechanism for natural user identification
and tracking. In Conference on Human Factors in Computing Systems (CHI2000), pages
1–6, 2000. ACM Press
[28] RODIN - Rigorous Open Development Environment for Complex Systems.
http://rodin.cs.ncl.ac.uk/
A. Matoussi, F. Gervais, R. Laleau. An Event-B formalization of KAOS goal refinement patterns 33
Appendix A. Another trace semantics to discover proof obligations related to the AND refinement
The proposed semantics consists to describe all the possible behaviors in terms of trace as
illustrated in the following diagram. Each behavior is described in the form of sequence of
events like the milestone refinement pattern.
Consequently, we have the following ”forward simulation” proof to perform in order to
ensure that any trace of a refined model must also be a trace of the abstraction:
(EvG1; EvG2) or (EvG2; EvG1)
For that, we extend the classical set theoretic representation (see Figure ??) as follows:
EvG1|||EvG2
∆= {w 7→ w
′
| ∃ v . I(v) ∧ J(v, w) ∧ (G1-Guard(w) ∧ G2-Guard(w))∧
((G1-Guard(w) ⇒ ∃ t . G1-PostCond(w, t)∧G2-Guard(t) ∧ G2-PostCond(t, w′
))∨
(G2-Guard(w) ⇒ ∃ t . G2-PostCond(w, t)∧G1-Guard(t) ∧ G1-PostCond(t, w′
)))}
Consequently, the two parts of the forward simulation condition can be stated as follows:
r
−1
; (EvG1|||EvG2) ∆= {v 7→ w
′
| ∃ w . I(v) ∧ J(v, w) ∧ (G1-Guard(w) ∧ G2-Guard(w))∧
((G1-Guard(w) ⇒ ∃ t . G1-PostCond(w, t)∧G2-Guard(t) ∧ G2-PostCond(t, w′
))∨
(G2-Guard(w) ⇒ ∃ t . G2-PostCond(w, t)∧G1-Guard(t) ∧ G1-PostCond(t, w′
)))}
EvG; r
−1 ∆= {v 7→ w
′
| ∃ v
′
.I(v) ∧ G-Guard(v) ∧ G-PostCond(v, v′
) ∧ I(v
′
) ∧ J(v
′
, w′
) }
Consequently, the translation of the forward simulation condition, namely r
−1
; (EvG1|||EvG2) ⊆
EvG; r
−1
, requires to prove the following sequent:
I(v) ∧ J(v, w) ∧ (G1-Guard(w) ∧ G2-Guard(w))∧
((G1-Guard(w) ⇒ ∃ t . G1-PostCond(w, t)∧G2-Guard(t) ∧ G2-
PostCond(t, w′
))∨
(G2-Guard(w) ⇒ ∃ t . G2-PostCond(w, t)∧G1-Guard(t) ∧ G1-
PostCond(t, w′
)))}
⊢
I(v) ∧ G-Guard(v) ∧ ∃ v
′
.G-PostCond(v, v′
) ∧ I(v
′
) ∧ J(v
′
, w′
)
(A.I)
The application of the inference rule OR-L allows us to simplify the above sequent by obtaining the two following sequents :
A. Matoussi, F. Gervais, R. Laleau. An Event-B formalization of KAOS goal refinement patterns 34
I(v) ∧ J(v, w) ∧ (G1-Guard(w) ∧ G2-Guard(w))∧
(G1-Guard(w) ⇒ ∃ t . G1-PostCond(w, t) ∧ G2-Guard(t) ∧ G2-
PostCond(t, w′
))
⊢
I(v) ∧ G-Guard(v) ∧ ∃ v
′
.G-PostCond(v, v′
) ∧ I(v
′
) ∧ J(v
′
, w′
)
(A.I.1)
I(v) ∧ J(v, w) ∧ (G1-Guard(w) ∧ G2-Guard(w))∧
(G2-Guard(w) ⇒ ∃ t . G2-PostCond(w, t) ∧ G1-Guard(t) ∧ G1-
PostCond(t, w′
))
⊢
I(v) ∧ G-Guard(v) ∧ ∃ v
′
.G-PostCond(v, v′
) ∧ I(v
′
) ∧ J(v
′
, w′
)
(A.I.2)
For each sequent ( (A.I.1) ,(A.I.2) ), we apply (in order) the inference rules: IMP-L then
MON then XST-L8
in order to simplify the hypothesis part of the sequent and also removing
useless hypotheses as follows:
I(v) ∧ J(v, w)∧
(G1-Guard(w) ∧ G1-PostCond(w, t) ∧ G2-PostCond(t, w′
))
⊢
I(v) ∧ G-Guard(v) ∧ ∃ v
′
.G-PostCond(v, v′
) ∧ I(v
′
) ∧ J(v
′
, w′
)
(A.I.1’)
The definition of the AND goal refinement pattern given by [4] help us to detect the
sufficient proof obligations in order to prove the last sequent; i.e. we ensure that proving
these proof obligations is sufficient to prove the sequent:
G1-Guard ⇒ G-Guard (PO1)
(G1-PostCond ∧ G2-PostCond) ⇒ G-PostCond (PO2)
I(v) ∧ J(v, w)∧
(G2-Guard(w) ∧ G2-PostCond(w, t) ∧ G1-PostCond(t, w′
)) ⊢
I(v) ∧ G-Guard(v) ∧ ∃ v
′
.G-PostCond(v, v′
) ∧ I(v
′
) ∧ J(v
′
, w′
)
(A.I.2’)
Similarly, the second sequent comes down to prove the following sufficient proof obligations:
G2-Guard ⇒ G-Guard (PO3)
(G2-PostCond ∧ G1-PostCond) ⇒ G-PostCond (PO2)
8
allows us to replace an existential assumption by one without the existential quantifier. This can only be
done however if the quantified variable is not free in the set of other assumptions and in the goal (see [3]).